\documentclass[journal]{IEEEtran}
\hyphenation{op-tical net-works semi-conduc-tor}
\usepackage[spanish]{babel}
\usepackage{graphicx}
\usepackage{amsmath}
\usepackage{hyperref}
\usepackage{float}

\begin{document}
\title{DISEÑO Y ANÁLISIS DE UN VEHÍCULO RC AUTÓNOMO - ENTREGABLE 1}
\author{Andrés~Botero, Jade~Moreno, Sofía~Reina and~Vivi~Delgado}% <-this % stops a space


% The paper headers
\markboth{Universidad de San Buenaventura - Ingeniería Multimedia }%
{Shell \MakeLowercase{\textit{Apellido, Nombre}}: Titulo del informe}
% The only time the second header will appear is for the odd numbered pages
% after the title page when using the twoside option.

% make the title area
\maketitle

\begin{abstract}
Se muestra el diseño electromecánico y el análisis cuantitativo de un vehículo a control remoto inspirado en el personaje Rayo McQueen, construido sobre un chasis de traccióna a cuatro ruedas (4WD) con cuatro motorreductores DC, dos controladores de puente H TB6612FNG y un microcontrolador ESP32. Se especifican el chasis, la tracción, los motores y los puentes H tipo MOSFET, la arquitectura de alimentación aislada mediante un regulador LM2596 para la potencia y un módulo dual para la lógica, el cableado y el firmware en C/C++ de tipo avance-giro-detención, sin control continuo ni sensores. Se explica el análisis cuantitativo de fuerza, torque, velocidad, relación de transmisión, Ley de Ohm y la corriente de arranque, contrastando los cálculos teóricos de requerimientos mecánicos con la validación experimental del prototipo funcional.
\end{abstract}

% Note that keywords are not normally used for peerreview papers.
\begin{IEEEkeywords}
Microcontrolador ESP32, Ley de Ohm, Motorreductor DC, Puente H MOSFET, Relación de transmisión, Torque, Vehículo RC.
\end{IEEEkeywords}



\section{Introducción}

\IEEEPARstart{L}{os} vehículos a control remoto (RC) son una plataforma didáctica habitual para integrar conceptos de mecánica, electrónica de potencia y programación embebida en un solo sistema físico. Este informe documenta la etapa de diseño electromecánico y firmware base de un vehículo RC con carrocería inspirada en el personaje Rayo McQueen (Cars, Pixar), manteniendo el enfoque centrado en los aspectos de ingeniería del prototipo como lo es: selección del chasis y tracción integral a cuatro ruedas (4WD), dimensionamiento de motores y transistores de alta eficiencia, arquitectura de alimentación, cableado y una rutina de firmware mínima gestionada por un microcontrolador ESP32.

\section{Objetivos}

\subsection{Objetivo general}
Desarrollar un vehículo a control remoto aplicando los conceptos básicos de mecánica, electrónica y programación.

\subsection{Objetivos específicos}
\begin{itemize}
    \item Construir un prototipo funcional capaz de ejecutar una secuencia autónoma simple (avanzar, girar, detenerse) sin retroalimentación de sensores ni de lazo de control continuo.
    \item Realizar el análisis cuantitativo completo de las magnitudes mecánicas y eléctricas involucradas.
    \item Implementar un firmware estructurado en C/C++ que gestione el control de los motores mediante PWM, garantizando la ejecución correcta de la secuencia de movimiento.
\end{itemize}

\section{Marco Teórico}
%Las variables van en cursiva y se realiza escribiendolas entre $$
\subsection{Torque y fuerza en motores DC}

Un motor DC con escobillas produce un torque proporcional a la corriente de armadura:

\begin{equation}
\tau = K_t \cdot I_a
\end{equation}

donde $K_t$ es la constante del torque del motor(N·m/A), $I_a$ es la corriente de armadura, según el modelo clásico de máquinas de corriente continua. \cite{chapman} El torque disponible en la rueda depende además de la relación de transmisión de la caja reductora integrada al motor.


\subsection{Relación de transmisión}
La relación de transmisión N (o"gear ratio") relaciona la velocidad agular del eje del motor con la del eje de salida, como lo afirma la teoría de trenes de engranajes \cite{norton}:  

\begin{equation}
N = \frac{\omega_{motor}}{\omega_{salida}} = \frac{\tau_{salida}}{\tau_{motor}} \quad \text{(idealmente, sin pérdidas)}
\end{equation}

En la práctica existe una eficiencia mecánica $\eta<1$ por fricción en los engranajes, de modo que el torque de salida real es \cite{norton}:
\begin{equation}
\tau_{salida} = N \cdot \tau_{motor} \cdot \eta
\end{equation}

\subsection{Cinemática rueda--chasis}
La velocidad lineal del vehículo se relaciona con la velocidad angular de la rueda de radio $r$ mediante la cinemática rotacional clásica \cite{norton}:

\begin{equation}
v = \omega \cdot r, \qquad \omega \, [\text{rad/s}] = \frac{2\pi \cdot N_{rpm}}{60}
\end{equation}

y la fuerza de tracción disponible en el contacto rueda--suelo se obtiene del torque de salida \cite{chapman}, \cite{norton}:

\begin{equation}
F = \frac{\tau_{salida}}{r}
\end{equation}

\subsection{Relación velocidad-torque (curva característica)}
Un motor DC de imán permanente presenta una relación lineal entre el torque y la velocidad angular, entre los puntos extremos de rotor bloqueado ($\tau_{arranque}$, $\omega=0$) y de vacío ($\tau=0$, $\omega_{\text{vacío}}$), tal como describe la curva característica par-velocidad de las máquinas DC \cite{chapman}:

\begin{equation}
\tau(\omega) = \tau_{arranque} \cdot \left(1 - \frac{\omega}{\omega_{\text{vacío}}}\right)
\end{equation}

La potencia mecánica entregada es $P = \tau \cdot \omega$, la cual --bajo este modelo lineal-- alcanza su máximo teórico en $\omega = \omega_{\text{vacío}}/2$, es decir con $\tau = \tau_{arranque}/2$ \cite{chapman}.

\subsection{Ley de Ohm, corriente de arranque y de régimen}
La corriente que consume el motor está gobernada por la Ley de Ohm aplicada al circuito de armadura, considerando la fuerza contraelectromotriz (fcem) generada por la rotación \cite{dorf}, \cite{nilsson}:

\begin{equation}
I = \frac{V - \text{fcem}}{R_a}
\end{equation}

Al momento del arranque ($\omega=0$) la fcem es nula, por lo que la corriente es máxima (corriente de arranque o de rotor bloqueado). A medida que el motor acelera, la fcem crece y la corriente disminuye hasta un valor de régimen (corriente en vacío o en carga estable), sensiblemente menor que la de arranque, comportamiento típico de circuitos con elementos inductivo--resistivos bajo Ley de Ohm generalizada \cite{dorf} \cite{nilsson}. Esta diferencia es crítica para dimensionar el puente H, que debe soportar el pico de arranque sin entrar en protección térmica.
 
\subsection{Puente H MOSFET}
El puente H permite invertir la polaridad aplicada al motor mediante cuatro interruptores de potencia organizados en dos ramas. A diferencia de las topologías basadas en transistores bipolares (BJT) tradicionales, los controladores modernos como el TB6612FNG emplean transistores MOSFET \cite{tb6612fng}. Esta arquitectura presenta una resistencia de encendido ($R_{DS(on)}$) muy baja, típicamente de 0.5 $\Omega$, lo que minimiza la disipación de calor y reduce drásticamente la caída de tensión interna respecto al voltaje de alimentación \cite{tb6612fng}. Su dimensionamiento requiere que la corriente máxima de arranque de la carga no supere la corriente pico soportada por el dispositivo, garantizando una operación térmica estable.

\subsection{Redes inalámbricas de área local: Wi-Fi (IEEE 802.11)}

Wi-Fi es el nombre comercial de la familia de estándares IEEE 802.11 para redes de área local inalámbrica \cite{ieee80211}. Define cómo se comunican los dispositivos por radio en bandas de uso libre, principalmente la de 2.4 GHz, que es la que emplean los módulos de este proyecto. Una red se identifica con un nombre (SSID), y cualquier equipo que lo conozca y disponga de la contraseña puede unirse a ella.

Existen dos formas de organizar la red:

\begin{itemize}
\item \textbf{Modo infraestructura:} el dispositivo se conecta a un punto de acceso ya existente, como el \emph{router} de una casa.
\item \textbf{Modo punto de acceso (AP o \emph{SoftAP}):} el propio dispositivo crea la red y actúa como punto de acceso. El vehículo queda así como el \emph{router} de su propia red y entrega automáticamente una dirección a cada equipo que se conecta, de modo que un teléfono o un computador se enlazan directamente con él sin depender de infraestructura externa.
\end{itemize}

Se eligió el modo punto de acceso porque en campo resulta más práctico que el propio vehículo cree la red, y porque toda la comunicación ocurre entre el cliente y la placa, con la red aislada de Internet. La red se protege con una contraseña (WPA2) y el acceso se controla desde la propia placa.

\subsection{Modelo TCP/IP}

Las redes funcionan por capas, y TCP/IP es el conjunto de protocolos que hace posible la comunicación en Internet y en casi cualquier red local actual. La primera pieza es \textbf{IP} (\emph{Internet Protocol}), que se encarga del direccionamiento: cada dispositivo tiene una dirección IP, y IP se ocupa de que la información llegue al destino correcto. Por encima de IP hay dos maneras de enviar los datos:

\begin{itemize}
\item \textbf{TCP} (\emph{Transmission Control Protocol}) garantiza que todo llegue: abre una conexión entre los dos extremos y se asegura de que los datos arriben completos y en orden, reenviando lo que se pierda.
\item \textbf{UDP} (\emph{User Datagram Protocol}) es más simple: envía cada mensaje por su cuenta, sin abrir antes una conexión ni confirmar que llegó, de modo que puede perderse o llegar desordenado \cite{rfc768}.
\end{itemize}

Para la transmisión en tiempo real, UDP tiene una ventaja clara. Al no esperar confirmaciones ni reenviar lo perdido, mantiene el retardo bajo y constante; TCP, en cambio, puede frenar todo lo que viene detrás de un paquete perdido hasta que lo recupera, de modo que la imagen se queda un instante congelada. En video o en datos en vivo esto importa mucho: un fotograma que llega tarde ya pierde su utilidad, así que lo práctico es descartarlo y dejar paso al siguiente.

En este proyecto, sin embargo, se optó por TCP por comodidad: HTTP ---el protocolo que se usa para la página y los comandos--- funciona sobre TCP, lo que permite emplear un servidor web estándar y un cliente tan simple como un navegador, a cambio de un retardo algo mayor. UDP queda como una mejora natural si más adelante se quiere reducir el retardo del video. Como opción a futuro, \textbf{WebRTC} (\emph{Web Real-Time Communication}) es el estándar pensado justo para eso: envía audio y video en tiempo real sobre UDP, ya incorpora su propio cifrado y control del tráfico, y los navegadores lo soportan de forma nativa \cite{rfc8825}. Adoptarlo reduciría notablemente el retardo del video frente a MJPEG sobre HTTP, a costa de un servidor más complejo.

En resumen, sobre esta base funcionan los protocolos de aplicación; en este proyecto, HTTP. Cuando un cliente pide una página o envía un comando, la placa se encarga de que los datos viajen y el protocolo de aplicación solo interpreta el contenido, sin que el usuario tenga que preocuparse por direcciones ni por datos perdidos.

\subsection{Protocolo de aplicación HTTP y estilo REST}

HTTP (\emph{Hypertext Transfer Protocol}) es el protocolo que usan las páginas web: el cliente pide algo y el servidor responde \cite{rfc9110}. Trabaja \textbf{sin estado}: cada petición se responde independientemente, así que el cliente incluye en ella todo lo necesario. Una petición indica qué se quiere hacer (el método) y sobre qué se quiere actuar (una dirección), y el servidor contesta con un resultado y su contenido.

Los métodos habituales son \texttt{GET}, para pedir información sin cambiarla, y \texttt{POST}, para ejecutar una acción o enviar datos. Esta separación tiene una razón práctica: las acciones con efecto real ---reproducir un sonido, mover los motores--- se envían con \texttt{POST}, de modo que ningún intermediario pueda repetirlas sin querer a partir de una copia guardada. La respuesta incluye un código que resume el resultado (todo bien, error del cliente o error del servidor) y un tipo de contenido que dice qué se está enviando: HTML para la página de control, JSON para los datos de estado o JPEG para una fotografía de la cámara.

El estilo REST organiza la aplicación como un conjunto de direcciones que se consultan y se modifican con esos mismos métodos \cite{fielding}. La placa expone así direcciones como \texttt{/status}, \texttt{/play}, \texttt{/control}, \texttt{/drive} y \texttt{/stop}, que el cliente consulta por HTTP.

Por último, el video en vivo se sirve como \emph{MJPEG}: HTTP deja abierta una única respuesta por la que van llegando los fotogramas uno tras otro \cite{rfc2046}. Como esa respuesta queda ocupada mientras alguien mira el video, se atiende en su propio puerto para no bloquear el resto de la aplicación.

\subsection{Conexiones serie: UART, I\textsuperscript{2}C e I\textsuperscript{2}S}

Dentro del vehículo, los periféricos se comunican por cables, y cada uno usa un tipo de conexión distinto según cuántos cables necesita y a qué velocidad va:

\begin{itemize}
\item \textbf{UART} (\emph{Universal Asynchronous Receiver--Transmitter}) es una conexión directa entre dos puntos en la que ambos acuerdan de antemano la velocidad, sin compartir una señal de reloj. Se emplea entre el cerebro y el cuerpo.
\item \textbf{I\textsuperscript{2}C} (bus de dos cables) permite conectar varios dispositivos compartiendo las mismas dos líneas, y cada uno responde a una dirección propia \cite{i2c}. Se usa para el sensor de inercia, y en este proyecto se comprobó que solo resulta fiable a la velocidad estándar de 100 kHz.
\item \textbf{I\textsuperscript{2}S} es la conexión de audio de tres cables ---uno de reloj, uno de sincronía y uno de datos--- que lleva el sonido hasta el convertidor de digital a analógico \cite{i2s}. La velocidad del reloj depende de la calidad del audio y de cuántos canales se transmitan:
\end{itemize}

\begin{equation}
f_{BCK} = f_s \cdot N_{canales} \cdot N_{bits}
\end{equation}

\noindent por ejemplo, a $f_s = 44.1$ kHz en estéreo de 16 bits resulta $f_{BCK} = 44\,100 \cdot 2 \cdot 16 = 1.4112$ MHz.

\section{Diseño y Metodología}
 
\subsection{Chasis y tracción}
Se emplea un chasis rígido de tracción a las cuatro ruedas (4WD) con cuatro motorreductores DC que accionan las ruedas de forma independiente, acoplados eléctricamente en paralelo por lado (dos derechos, dos izquierdos) para lograr un giro diferencial. La carrocería estética (Rayo McQueen) se monta de forma no estructural para no restringir el acceso al cableado.
 
\subsection{Motores DC y caja reductora}
Se seleccionan motorreductores tipo ``TT'', habituales en kits educativos de robótica móvil, con caja reductora integrada de relación nominal 1:48 y tensión de operación 3--6 V. Estos valores corresponden a especificaciones típicas reportadas por múltiples fabricantes de este tipo de motorreductor y deben confirmarse contra la hoja de datos específica del lote adquirido.
 
\subsection{Puente H y dimensionamiento}
Dado que se manejan cuatro motores, se utilizan dos módulos TB6612FNG (un módulo para el flanco derecho y otro para el izquierdo). Según el fabricante, cada canal soporta 1.2 A continuos y hasta 3.2 A pico \cite{tb6612fng}. Al distribuir la carga en dos controladores, se garantiza que la corriente de arranque de los motores en paralelo no exceda los límites operativos ni active la protección térmica de los integrados.Esta capacidad de manejo de corriente eficiente es posible gracias a la arquitectura interna basada en MOSFETs de baja resistencia, ilustrada en el diagrama de bloques del controlador (ver Figura \ref{fig:tb6612fng}) \cite{tb6612fng}."

\begin{figure}[H]
    \centering
    \includegraphics[width=0.95\linewidth]{Block Diagram.png}
    \caption{Diagrama de bloques interno del controlador TB6612FNG.}
    \label{fig:tb6612fng}
\end{figure}

\subsection{Arquitectura de alimentación independiente}
Se diseñó una arquitectura de potencia aislada a partir de una única batería de polímero de litio (LiPo) de alto rendimiento, marca CNHL (grado dron), de 3 celdas (3S), con una tensión nominal de 11.1 V y una capacidad de 2200 mAh. La elección de una batería orientada a vehículos aéreos no tripulados garantiza una alta tasa de descarga continua, capaz de suministrar los picos de corriente demandados por los cuatro motorreductores durante el arranque sin sufrir caídas abruptas de tensión. A partir de esta fuente primaria, el sistema se ramifica compartiendo una tierra común:

\begin{itemize}
\item \textbf{Potencia de motores:} El voltaje de 11.1 V de la batería CNHL ingresa a un módulo regulador DC-DC buck LM2596 \cite{lm2596}, el cual reduce eficientemente la tensión a un nivel seguro (ej. 6 V) para la entrada de potencia ($VM$) de los módulos TB6612FNG y, en consecuencia, de los motores.El circuito reductor implementado sigue la topología de referencia del fabricante para garantizar una salida estable ante fluctuaciones de carga, tal como se detalla en la Figura \ref{fig:lm2596} \cite{lm2596}.".

\begin{figure}[H]
    \centering
    \includegraphics[width=1\linewidth]{Typical Application.png}
    \caption{Circuito típico de aplicación del regulador Buck LM2596.}
    \label{fig:lm2596}
\end{figure}


\item \textbf{Lógica/control:} En paralelo, la alimentación principal ingresa a un módulo regulador de protoboard dual que entrega 3.3 V estables para alimentar el microcontrolador ESP32 \cite{esp32} y los pines de referencia lógica ($VCC$) de los puentes H.
\end{itemize}
Esta separación por conmutación, sumada a la rigidez dieléctrica y alta capacidad de respuesta de la batería LiPo, evita que los transitorios mecánicos de los motores provoquen reinicios (brownouts) en la lógica del microcontrolador.
 
\subsection{Cableado}
El cableado de potencia se dimensiona en calibre suficiente para la corriente de arranque estimada y se mantiene físicamente separado del cableado de señal (ESP32 a TB6612FNG) para reducir el acoplamiento de ruido eléctrico durante la conmutación de los motores.

\subsection{Arquitectura de control distribuida: placas utilizadas y su función}

El Entregable 1 describía un único microcontrolador ESP32 clásico que generaba directamente las señales PWM y de dirección de los dos puentes H. Al incorporar la cámara, el audio, la unidad de medición inercial y la tira de LEDs, nos quedamos cortos de pines. La cámara ocupa de forma fija 14 pines, el audio (I\textsuperscript{2}S) necesita 3 (reloj, sincronía y datos); el sensor de inercia (I\textsuperscript{2}C) necesita 2 (SDA y SCL); y la tira de LEDs, 1. Eso suma 20 pines para los periféricos. Del lado de los motores, controlar los cuatro motorreductores con dos puentes dobles exige 12 líneas ---un PWM y dos de dirección por cada uno de los cuatro canales--- más 1 línea de \emph{standby} compartida, es decir 13 pines adicionales. En total son 33 pines, muy por encima de los GPIO que quedan libres en el ESP32-S3 una vez descontados los de la cámara, la memoria flash y la PSRAM. Por esta razón la arquitectura se dividió en dos microcontroladores con responsabilidades separadas:

\begin{itemize}
\item Un \textbf{cerebro} (ESP32-S3-CAM) que \emph{decide}: red Wi-Fi, captura y transmisión de cámara, reproducción de audio, lectura del sensor de movimiento y el control del vehículo (mezcla, rampas y frenado).
\item Un \textbf{cuerpo} (RP2040, Raspberry Pi Pico) que se limita a ejecutar: recibe cuatro intensidades de potencia por puerto serial y las traduce en PWM y líneas de dirección, aportando además la seguridad de un \emph{dead-man} (vigilante de comunicación) implementado por firmware.
\end{itemize}

La Tabla~\ref{tab:placas} resume las placas y el uso dado a cada una.

\begin{table}[H]
\centering
\caption{Placas utilizadas y su función dentro del sistema.}
\label{tab:placas}
\small
\begin{tabular}{|p{2.1cm}|p{1.5cm}|p{4.3cm}|}
\hline
\textbf{Placa} & \textbf{MCU} & \textbf{Función en el proyecto} \\
\hline
ESP32-S3-CAM (WROOM-1 N16R8) & ESP32-S3 & Cerebro: Wi-Fi (SoftAP), cámara, audio, sensor de movimiento y control del vehículo; puerto serial con el cuerpo. \\
\hline
Raspberry Pi Pico & RP2040 & Cuerpo: 4 PWM a 20 kHz + 8 líneas de dirección + \emph{STBY} para los dos ROB-14450; vigilante (\emph{dead-man}) de 100 ms. \\
\hline
ESP32 clásico (\texttt{esp32dev}) & ESP32 & Prototipo del Entregable 1; retirado del vehículo. Su lógica de mezcla se portó al cerebro. \\
\hline
\end{tabular}
\end{table}

El cerebro es un ESP32-S3 de dos núcleos con 16 MB de memoria flash y 8 MB de PSRAM, montado sobre una placa portadora de 40 pines (Freenove ESP32-S3-WROOM) a la que se conecta la cámara por un conector plano. La PSRAM es indispensable para guardar las imágenes de la cámara mientras se procesan, y los dos núcleos permiten repartir las tareas: la red y el audio por un lado, el control del movimiento por el otro. El cuerpo es un RP2040 (dos núcleos Cortex-M0+ a 133 MHz) elegido porque genera el PWM del propio chip, a cualquier frecuencia y sin cargar al cerebro, que solo le dedica los dos pines del puerto serial. El detalle de las conexiones y la asignación de pines se documenta en \texttt{docs/circuit-plan.md} y en el esquema eléctrico del repositorio.

\subsection{Dispositivos de entrada y salida}

Además de los puentes H y los motores, el vehículo integra cuatro dispositivos que amplían su percepción y su capacidad de interacción. La Tabla~\ref{tab:io} los resume y las subsecciones siguientes detallan cada uno.

\begin{table}[H]
\centering
\caption{Dispositivos de entrada y salida del vehículo.}
\label{tab:io}
\small
\begin{tabular}{|p{2.0cm}|p{1.9cm}|p{4.0cm}|}
\hline
\textbf{Dispositivo} & \textbf{Modelo} & \textbf{Cómo se conecta y para qué sirve} \\
\hline
Cámara (entrada) & OV3660 & Conexión paralela de la propia placa; guarda las imágenes JPEG en la PSRAM y las transmite en vivo. \\
\hline
Sensor de movimiento (entrada) & BMI160 (6 ejes) & I\textsuperscript{2}C a 100 kHz; orientación, inclinación y detección de golpes. \\
\hline
Audio (salida) & PCM5102A + amplificador clase D & I\textsuperscript{2}S hacia el convertidor de audio y el parlante. \\
\hline
Luces (salida) & Tira WS2812 y LED RGB integrado & RMT (generador de pulsos) en GPIO48; estado y alumbrado. \\
\hline
\end{tabular}
\end{table}

\subsubsection{Cámara (entrada)}
El sensor de imagen es un \textbf{OV3660} \cite{ov3660}. Se conecta a la placa por una línea paralela propia de la cámara, y la placa lo configura por una línea de control aparte. El sensor entrega directamente las imágenes ya comprimidas en JPEG, y el ESP32-S3 las va guardando en la PSRAM. Como viene montado girado sobre la placa, el firmware invierte la imagen al encenderla para que se vea derecha. El resultado se ofrece de dos maneras: como video en vivo (\texttt{/stream}, MJPEG) y como fotografía única (\texttt{/capture}). Su función es dar al operador una vista de lo que hay delante y servir de base para desarrollos futuros de visión.

\subsubsection{Audio (salida)}
El audio se arma en cadena: un convertidor de digital a analógico \textbf{PCM5102A} \cite{pcm5102a} recibe el sonido ya digitalizado por la conexión I\textsuperscript{2}S y lo entrega como una señal de nivel de línea. Por eso se añade un \textbf{amplificador clase D} \cite{pam8302a} y, a su salida, el parlante. El sonido se guarda temporalmente en la memoria interna del microcontrolador para no competir con la PSRAM que usa la cámara, y se reproduce en una tarea aparte para no bloquear el servidor web. La reproducción es bajo demanda: el vehículo solo suena cuando se le pide. Su función es reproducir efectos sonoros y avisos.

\subsubsection{Luces (salida)}
La iluminación usa una tira de LEDs de color controlables \textbf{WS2812} de 5 V \cite{ws2812}, montada en paralelo con el LED RGB que ya trae la placa: ambos reciben los datos por el mismo pin, el GPIO48. Estas tiras requieren una señal de un solo cable en la que la duración de cada pulso debe ser muy exacta, por lo que se genera con el periférico \textbf{RMT} del ESP32-S3 para sincronizarse de manera exacta. Su función principal es dar personalidad al carro, y pueden alumbrar para permitir que la cámara vea en la oscuridad.

\subsubsection{Sensor de movimiento (entrada)}
El sensor de movimiento es un \textbf{BMI160} de 6 ejes ---acelerómetro y giroscopio--- \cite{bmi160}. Se comunica por I\textsuperscript{2}C. Toma 200 muestras por segundo, con un alcance de $\pm 8$~g en las aceleraciones y $\pm 2000$~ grados/segundo en el giro, para saber cómo está orientado e inclinado el vehículo. Estos datos permiten mostrar una interfaz de inclinación en la aplicación de Godot.

\subsection{Firmware (rutina mínima avance--giro--detención)}


El firmware es implementado en C/C++ para el microcontrolador ESP32 \cite{esp32}. Tiene una secuencia de arranque, donde los 4 motores se activan suavemente para probar que el PWM este funcionando correctamente y la direccion de rotacion sea la correcta. Después de esto se ejecuta la secuencia principal, donde hay varias tareas previamente implementadas:
\begin{itemize}
\item Movimiento lento hacia adelante y atrás.
\item Movimiento rápido hacia adelante y atrás.
\item Giro sobre el propio eje del vehiculo, en ambos sentidos.
\item Patron de estrella (movimiento en curva hacia el frente y hacia atrás). 

\end{itemize}
Para la compilación, gestión de dependencias y despliegue del código en el microcontrolador se empleó el ecosistema PlatformIO integrado en Visual Studio Code (VSCode) \cite{platformio}. Esta herramienta ofrece soporte nativo para el ESP32, permitiendo un control estricto de las librerías del framework y la automatización del proceso de carga mediante interfaz serial, garantizando un entorno de desarrollo reproducible y profesional frente a los entornos de desarrollo tradicionales de Arduino.

Los pines de control se asignan a puertos GPIO digitales aptos para  modulación por ancho de pulso (PWM). La selección de estos puertos se realizó respetando la distribución de periféricos del microcontrolador, como se muestra en el esquema de pines de la Figura \ref{fig:esp32} \cite{esp32}:

\begin{figure}[H]
    \centering
    \includegraphics[width=1\linewidth]{ESP32-Board-Pinout.png}
    \caption{Distribución de pines (Pinout) del microcontrolador ESP32 destacando puertos GPIO/PWM.}
    \label{fig:esp32}
\end{figure}

La seleccion de estos puertos se hizo tomando los pines mas cercanos a los modulos puente H:

\begin{verbatim}
struct Motor {
    int pin_pwm;
    int pin_dir_a;
    int pin_dir_b;
};

Motor motors[] = {
    { 23, 21, 22 },
    { 4,  18, 19 },
    { 14, 26, 27 },
    { 32, 25, 33 },
};
\end{verbatim}

La rutina ciclica del robot (loop) ejecuta una lista de tareas establecidas, donde cada una tiene una duracion, velocidad base, y velocidades para pares diferenciales (izquierda/derecha, frente/atras):

\begin{verbatim}
struct Task {
    int time_ms;
    int input_all;
    int input_front_rear; // added to front 
    wheels, subtracted from rear wheels
    int input_left_right; // added to left 
    wheels, subtracted from right wheels
};

Task tasks[] = {
    // slow fwd
    { 1000, 128, 0, 0},
    // slow back
    { 1000, -128, 0, 0},
...
    // turn in place left
    { 1000, 0, 0, 256 },
    // turn in place right
    { 1000, 0, 0, -256 },
...
}
\end{verbatim}
\textbf{El codigo completo se puede encontrar en el repositorio:} https://github.com/botero-dev/usbbog-realidad-mixta

\section{Resultados}

El diseño electromecánico detallado en la sección anterior se materializó en un prototipo funcional, cuya integración física se presenta en la Figura \ref{fig:ensamble_fisico}. En esta vista superior se evidencia la distribución del cableado de potencia y señal, la ubicación central del microcontrolador ESP32 y la disposición simétrica de los controladores de motor para la tracción 4WD.

\begin{figure}[H]
    \centering
    \includegraphics[width=0.8\linewidth]{ensamble_fisico.jpeg}
    \caption{Vista superior del ensamblaje electromecánico final del vehículo RC, mostrando la integración del chasis, motores, puentes H, regulador de voltaje y microcontrolador.}
    \label{fig:ensamble_fisico}
\end{figure}



\subsection{Análisis de Torque y Tracción}
La masa total medida del vehículo ensamblado (incluyendo chasis, motores, puente H, ESP32 y batería LiPo) es de $m = 0.63$ kg. Considerando el radio medido de las ruedas $r = 0.03375$ m (equivalente a un diámetro de $67.5$ mm) y un coeficiente de fricción estimado $\mu = 0.6$ para rodadura en superficie plana, la fuerza de tracción mínima requerida es:

$$F_{min} = m \cdot g \cdot \mu = 0.63 \cdot 9.81 \cdot 0.6 = 3.70 \text{ N}$$

El torque total necesario en el eje de las ruedas es $\tau_{req} = F_{min} \cdot r = 3.70 \text{ N} \cdot 0.03375 \text{ m} = 0.125 \text{ N}\cdot\text{m}$. Dado que el sistema 4WD distribuye la carga, cada uno de los cuatro motores TT proporciona un torque combinado de salida superior al requerido, garantizando el desplazamiento del vehículo.

\subsection{Análisis Eléctrico y Dimensionamiento}
Según las especificaciones técnicas, el motor TT estándar a 6 V presenta una corriente de arranque de aproximadamente $1.0 \text{ A}$ por unidad. Al estar acoplados en paralelo por lado, cada puente H TB6612FNG experimenta un pico de corriente de $2.0 \text{ A}$ en el arranque. Dado que el controlador TB6612FNG soporta hasta $3.2 \text{ A}$ pico por canal \cite{tb6612fng}, el sistema está correctamente dimensionado y opera dentro de un margen seguro frente a su límite térmico, validando la arquitectura de alimentación propuesta.

En un momento dañamos el regulador de voltaje al conectarle el voltaje de entrada con la polaridad invertida. Por esto tuvimos que buscar opciones para poder alimentar el carro robot. Una de las opciones fue usar un juego de 4 pilas AA en serie. En principio, esto parecía que iba a funcionar, ya que las ruedas giraban bien mientras programabamos el microcontrolador, y teníamos el carro en un pedestal en la mesa de trabajo. Pero en cuanto pusimos el carro sobre sus ruedas, el consumo de corriente de los motores reducía tanto el voltaje debido al peso de sus componentes, y así mismo interfiere la fricción entre la superficie y el material de la llanta del vehículo, provocando que se reiniciara el microcontrolador. Afortunadamente encontramos un par de baterías 18650, y de esta manera logramos alimentar el carro con un voltaje adecuado para los motores, alrededor de 7.4V. Con este par de baterias hicimos la prueba final de funcionamiento que se ve en el video.


\textbf{Video del vehículo en funcionamiento mostrando la ejecución completa de la rutina automática:} https://youtu.be/Z-YuQ6jdH08



% references section

% can use a bibliography generated by BibTeX as a .bbl file
% BibTeX documentation can be easily obtained at:
% http://mirror.ctan.org/biblio/bibtex/contrib/doc/
% The IEEEtran BibTeX style support page is at:
% http://www.michaelshell.org/tex/ieeetran/bibtex/
%\bibliographystyle{IEEEtran}
% argument is your BibTeX string definitions and bibliography database(s)
%\bibliography{IEEEabrv,../bib/paper}
%
% <OR> manually copy in the resultant .bbl file
% set second argument of \begin to the number of references
% (used to reserve space for the reference number labels box)
\begin{thebibliography}{00}

\bibitem{chapman} S. J. Chapman, \emph{Electric Machinery Fundamentals}, 5th ed. New York, NY, USA: McGraw-Hill, 2012.

\bibitem{dorf} R. C. Dorf and J. A. Svoboda, \emph{Introduction to Electric Circuits}, 9th ed. Hoboken, NJ, USA: Wiley, 2013.

\bibitem{tb6612fng} Toshiba Corporation, ``TB6612FNG: Driver IC for Dual DC motor,'' Datasheet, Oct. 2014. [Online]. Available: https://toshiba.semicon-storage.com/

\bibitem{esp32} Espressif Systems, ``ESP32 Series Datasheet,'' v5.0, 2023. [Online]. Available: https://www.espressif.com/

\bibitem{norton} R. L. Norton, \emph{Design of Machinery}, 5th ed. New York, NY, USA: McGraw-Hill, 2012.

\bibitem{nilsson} J. W. Nilsson and S. A. Riedel, \emph{Electric Circuits}, 10th ed. Boston, MA, USA: Pearson, 2015.

\bibitem{lm2596} Texas Instruments, ``LM2596 SIMPLE SWITCHER Power Converter 150-kHz 3-A Step-Down Voltage Regulator,'' Datasheet, Rev. G, 2020. [Online]. Available: https://www.ti.com/product/LM2596

\bibitem{platformio} PlatformIO Labs, ``PlatformIO: Your permanent home for open source software debuggers and toolchains,'' Online Documentation, 2026. [Online]. Available: https://platformio.org/

\bibitem{ieee80211} IEEE, ``IEEE Standard for Information Technology---Telecommunications and Information Exchange between Systems---Local and Metropolitan Area Networks---Specific Requirements---Part 11: Wireless LAN Medium Access Control (MAC) and Physical Layer (PHY) Specifications,'' IEEE Std 802.11-2020, 2021.

\bibitem{rfc9110} R. Fielding, M. Nottingham, and J. Reschke, ``HTTP Semantics,'' RFC 9110, IETF, Jun. 2022. [Online]. Available: https://www.rfc-editor.org/rfc/rfc9110

\bibitem{rfc2046} N. Freed and N. Borenstein, ``Multipurpose Internet Mail Extensions (MIME) Part Two: Media Types,'' RFC 2046, IETF, Nov. 1996. [Online]. Available: https://www.rfc-editor.org/rfc/rfc2046

\bibitem{rfc768} J. Postel, ``User Datagram Protocol,'' RFC 768, IETF, Aug. 1980. [Online]. Available: https://www.rfc-editor.org/rfc/rfc768

\bibitem{rfc8825} C. Jennings, H. Bostr{\"o}m, and J.-I. Bruaroey, ``WebRTC: Real-Time Communication in Browsers,'' RFC 8825, IETF, Jan. 2021. [Online]. Available: https://www.rfc-editor.org/rfc/rfc8825

\bibitem{fielding} R. T. Fielding, ``Architectural Styles and the Design of Network-based Software Architectures,'' Ph.D. dissertation, Univ. of California, Irvine, Irvine, CA, USA, 2000.

\bibitem{i2c} NXP Semiconductors, ``UM10204: I\textsuperscript{2}C-bus specification and user manual,'' Rev. 7.0, 2021. [Online]. Available: https://www.nxp.com/docs/en/user-guide/UM10204.pdf

\bibitem{i2s} Philips Semiconductors, ``I\textsuperscript{2}S bus specification,'' Rev. Feb. 1986.

\bibitem{esp32s3} Espressif Systems, ``ESP32-S3 Series Datasheet,'' Espressif Systems, 2024. [Online]. Available: https://www.espressif.com/

\bibitem{rp2040} Raspberry Pi Ltd, ``RP2040 Datasheet: A microcontroller by Raspberry Pi,'' Raspberry Pi Ltd, 2024. [Online]. Available: https://datasheets.raspberrypi.com/rp2040/rp2040-datasheet.pdf

\bibitem{ov3660} OmniVision Technologies, ``OV3660 1/5-inch 3-megapixel CMOS image sensor,'' Datasheet. [Online]. Available: https://www.ovt.com/

\bibitem{bmi160} Bosch Sensortec, ``BMI160: Small, low power inertial measurement unit,'' Datasheet, Rev. 1.0. [Online]. Available: https://www.bosch-sensortec.com/

\bibitem{pcm5102a} Texas Instruments, ``PCM5102A: 2-V\textsubscript{RMS} Audio Stereo DAC with 32-bit, 384-kHz PCM Interface,'' Datasheet. [Online]. Available: https://www.ti.com/product/PCM5102A

\bibitem{pam8302a} Diodes Incorporated, ``PAM8302A: 2.5 W Filterless Mono Class-D Audio Amplifier,'' Datasheet. [Online]. Available: https://www.diodes.com/

\bibitem{ws2812} Worldsemi, ``WS2812B: Intelligent control LED integrated light source,'' Datasheet. [Online]. Available: https://www.worldsemi.com/

\end{thebibliography}

% that's all folks
\end{document}


